\documentclass[11pt]{article}

\usepackage[margin=1in]{geometry}
\usepackage{amsmath,amssymb,amsthm,mathtools}
\usepackage[T1]{fontenc}
\usepackage{lmodern}
\usepackage{microtype}

\newcommand{\OO}{\mathcal O}
\newcommand{\Bl}{\operatorname{Bl}}

\title{Natural-language proof of \texttt{IdealSheafData.exists\_blowup\_mul\_iso}}
\author{}
\date{}

\begin{document}

\maketitle

\section*{Statement}

Let \(X\) and \(X_1\) be schemes, with \(X\) quasi-compact and quasi-separated.
Let \(I\) be a finite-type ideal sheaf on \(X\), let
\[
  \pi : X_1 \longrightarrow X,
\]
and suppose that
\[
  g : X_1 \xrightarrow{\sim} \Bl_I(X)
\]
is an isomorphism satisfying
\[
  g \circ I.\mathrm{blowupMap} \text{ in Lean notation, equivalently }
  g.\mathrm{hom} \gg I.\mathrm{blowupMap}=\pi.
\]
Let \(L\) be another finite-type ideal sheaf on \(X\).

Then there is an isomorphism
\[
  e :
  \Bl_{I L}(X)
  \xrightarrow{\sim}
  \Bl_{\pi^{-1}L}(X_1)
\]
such that, as morphisms to \(X\),
\[
  e.\mathrm{hom}
  \gg
  (L.\mathrm{comap}\,\pi).\mathrm{blowupMap}
  \gg
  \pi
  =
  (I*L).\mathrm{blowupMap}.
\]

Here \(\pi^{-1}L\) means the inverse-image ideal
\(L.\mathrm{comap}\,\pi\).

\section*{Proof}

Write
\[
  b_I : \Bl_I(X)\longrightarrow X
\]
for the blow-up of \(I\), and use the given isomorphism
\[
  g:X_1\xrightarrow{\sim}\Bl_I(X)
\]
to regard \(\pi:X_1\to X\) as a chosen model of this blow-up.  Thus
\[
  \pi=g\circ b_I
\]
in ordinary geometric notation, or
\(g.\mathrm{hom}\gg I.\mathrm{blowupMap}=\pi\) in Lean's composition convention.

Now let
\[
  q:Y\longrightarrow X_1
\]
be the blow-up of the inverse-image ideal
\[
  L_1:=L.\mathrm{comap}\,\pi.
\]
Thus
\[
  Y=(L.\mathrm{comap}\,\pi).\mathrm{blowup}.
\]
We shall prove that the composite
\[
  \rho:=q\circ\pi:Y\longrightarrow X
\]
is, up to unique isomorphism over \(X\), the blow-up of the product ideal
\(I L\).

\subsection*{1. The product ideal becomes invertible on the iterated blow-up}

Because \(\pi:X_1\to X\) is the blow-up of \(I\), the inverse image of \(I\)
on \(X_1\),
\[
  I_1:=I.\mathrm{comap}\,\pi,
\]
is an invertible ideal sheaf.  Equivalently, it is the ideal sheaf of the
exceptional effective Cartier divisor of the first blow-up.

Pulling an invertible ideal back along an arbitrary morphism preserves
invertibility.  Hence
\[
  q^{-1}I_1\cdot\OO_Y
\]
is an invertible ideal on \(Y\).

By construction of \(q\), the inverse image of \(L_1\) on \(Y\),
\[
  q^{-1}L_1\cdot\OO_Y,
\]
is also an invertible ideal: it is the exceptional ideal of the second
blow-up.

The inverse image of a product of quasi-coherent ideals is the product of
their inverse images.  Therefore the pullback of \(I L\) along
\(\rho=\pi\circ q\) is
\[
\begin{aligned}
  \rho^{-1}(I L)\cdot\OO_Y
  &=
  \bigl(\rho^{-1}I\cdot\OO_Y\bigr)
  \bigl(\rho^{-1}L\cdot\OO_Y\bigr)\\
  &=
  \bigl(q^{-1}I_1\cdot\OO_Y\bigr)
  \bigl(q^{-1}L_1\cdot\OO_Y\bigr).
\end{aligned}
\]
Both factors are invertible ideals, hence so is their product.

Consequently the morphism
\[
  \rho:Y\longrightarrow X
\]
has the defining property required by the universal property of the
blow-up of \(I L\): the inverse image of \(I L\) is invertible.  Hence there
is a unique morphism
\[
  \alpha:
  Y\longrightarrow \Bl_{I L}(X)
\]
over \(X\).  In particular,
\[
  \alpha\circ b_{I L}=\rho
\]
with the appropriate convention for composition; in Lean notation this
says
\[
  \alpha \gg (I*L).\mathrm{blowupMap}
  =
  q\gg\pi.
\]

\subsection*{2. Recovering the first blow-up from the product blow-up}

Conversely, let
\[
  c:Z\longrightarrow X
\]
denote the blow-up of the product ideal \(I L\):
\[
  Z=(I*L).\mathrm{blowup},
  \qquad
  c=(I*L).\mathrm{blowupMap}.
\]
By definition of the blow-up,
\[
  c^{-1}(I L)\cdot\OO_Z
\]
is an invertible ideal sheaf on \(Z\).

We claim that
\[
  c^{-1}I\cdot\OO_Z
\]
is itself invertible.  This is the standard ``factor of an invertible
product ideal'' lemma used in the proof of Stacks Project, Tag 080A.
Indeed, locally on \(Z\), if the product
\[
  \bigl(c^{-1}I\cdot\OO_Z\bigr)
  \bigl(c^{-1}L\cdot\OO_Z\bigr)
\]
is an invertible ideal, then the first factor defines an effective Cartier
divisor; equivalently, it is an invertible ideal.  Geometrically, the
closed subscheme defined by the first factor is an effective Cartier
divisor because its sum with the closed subscheme defined by the second
factor is the effective Cartier divisor defined by the product.

Thus \(c^{-1}I\cdot\OO_Z\) is invertible.  Applying the universal property
of the blow-up of \(I\), there is a unique morphism
\[
  \beta_0:Z\longrightarrow \Bl_I(X)
\]
over \(X\).

Using the given isomorphism
\[
  g:X_1\xrightarrow{\sim}\Bl_I(X),
\]
define
\[
  \beta:=\beta_0\circ g^{-1}:Z\longrightarrow X_1.
\]
Then \(\beta\) lies over \(X\):
\[
  \pi\circ\beta=c.
\]
Equivalently, in Lean notation,
\[
  \beta\gg\pi=(I*L).\mathrm{blowupMap}.
\]

\subsection*{3. The pullback of \(L\) is invertible after the first factorization}

It remains to show that \(\beta:Z\to X_1\) factors through the blow-up of
\(L_1=L.\mathrm{comap}\,\pi\).

The inverse image of \(L_1\) along \(\beta\) is exactly the inverse image of
\(L\) along \(c=\pi\circ\beta\):
\[
  \beta^{-1}L_1\cdot\OO_Z
  =
  c^{-1}L\cdot\OO_Z.
\]
We therefore need to show that \(c^{-1}L\cdot\OO_Z\) is invertible.

Again,
\[
  \bigl(c^{-1}I\cdot\OO_Z\bigr)
  \bigl(c^{-1}L\cdot\OO_Z\bigr)
  =
  c^{-1}(I L)\cdot\OO_Z
\]
is invertible.  We have already shown that the first factor
\(c^{-1}I\cdot\OO_Z\) is invertible.  Tensoring, or equivalently multiplying
by its inverse line bundle, shows that the second factor
\(c^{-1}L\cdot\OO_Z\) is invertible as well.

Hence the universal property of the blow-up
\[
  q:Y=\Bl_{L_1}(X_1)\longrightarrow X_1
\]
gives a unique morphism
\[
  \gamma:Z\longrightarrow Y
\]
over \(X_1\).  Therefore
\[
  q\circ\gamma=\beta.
\]
Composing with \(\pi\), we obtain
\[
  \rho\circ\gamma
  =
  \pi\circ q\circ\gamma
  =
  \pi\circ\beta
  =
  c.
\]
Thus \(\gamma\) is also a morphism over \(X\).

\subsection*{4. The two morphisms are inverse}

We have constructed morphisms
\[
  \alpha:Y\longrightarrow Z
  \qquad\text{and}\qquad
  \gamma:Z\longrightarrow Y,
\]
where
\[
  Y=\Bl_{L.\mathrm{comap}\,\pi}(X_1),
  \qquad
  Z=\Bl_{I L}(X).
\]

Consider \(\alpha\circ\gamma:Z\to Z\).  Both this composite and
\(\operatorname{id}_Z\) are morphisms over \(X\).  Moreover, \(Z\) is the
blow-up of \(I L\), hence it is final among \(X\)-schemes on which the
inverse image of \(I L\) is invertible.  Since the inverse image of \(I L\)
on \(Z\) is invertible, there is only one morphism from \(Z\) to the
blow-up \(Z\) over \(X\).  Therefore
\[
  \alpha\circ\gamma=\operatorname{id}_Z.
\]

For the other composite, \(\gamma\circ\alpha:Y\to Y\), we use the second
blow-up universal property.  Both \(\gamma\circ\alpha\) and
\(\operatorname{id}_Y\) are morphisms over \(X_1\): indeed \(\gamma\) was
constructed over \(X_1\), and \(\alpha\) respects the factorization through
the first blow-up by uniqueness in the universal property of
\(\Bl_I(X)\).  Since \(Y\) is the blow-up of \(L_1\) and the inverse image
of \(L_1\) on \(Y\) is invertible, uniqueness gives
\[
  \gamma\circ\alpha=\operatorname{id}_Y.
\]

Thus \(\gamma\) and \(\alpha\) are inverse isomorphisms.  Taking
\[
  e:=\gamma:
  (I*L).\mathrm{blowup}
  \xrightarrow{\sim}
  (L.\mathrm{comap}\,\pi).\mathrm{blowup},
\]
we obtain the required isomorphism.

Finally, because \(\gamma\) was constructed over \(X_1\) and
\(\beta\) was constructed over \(X\), we have
\[
  e.\mathrm{hom}
  \gg
  (L.\mathrm{comap}\,\pi).\mathrm{blowupMap}
  \gg
  \pi
  =
  (I*L).\mathrm{blowupMap},
\]
which is precisely the compatibility required in the theorem.

\hfill\(\Box\)

\section*{Formalization-oriented summary}

For a Lean formalization, the proof naturally separates into the following
universal-property steps.

\begin{enumerate}
  \item On
  \[
    Y=(L.\mathrm{comap}\,\pi).\mathrm{blowup},
  \]
  show that the pullback of \(I*L\) along
  \[
    Y\to X_1\overset{\pi}{\to}X
  \]
  is invertible.  The \(I\)-factor is invertible because \(X_1\) is
  isomorphic over \(X\) to the blow-up of \(I\); the \(L\)-factor is
  invertible by the defining property of the second blow-up.

  \item Apply the universal property of \((I*L).\mathrm{blowup}\) to obtain
  \[
    Y\to (I*L).\mathrm{blowup}.
  \]

  \item On
  \[
    Z=(I*L).\mathrm{blowup},
  \]
  use the standard lemma that if the pullback of \(I*L\) is invertible,
  then the pullback of \(I\) is invertible.  Apply the universal property
  of \(I.\mathrm{blowup}\), and transport the resulting morphism through
  \(g^{-1}\), obtaining
  \[
    Z\to X_1.
  \]

  \item Show that the pullback of \(L.\mathrm{comap}\,\pi\) along this map
  is invertible.  Since the pullback of \(I*L\) and the pullback of \(I\)
  are invertible, so is the pullback of \(L\).

  \item Apply the universal property of
  \((L.\mathrm{comap}\,\pi).\mathrm{blowup}\) to obtain
  \[
    Z\to Y.
  \]

  \item Prove the two composites are identities by uniqueness in the two
  relevant blow-up universal properties.

  \item The required equality of structure morphisms is built into the
  construction of the map \(Z\to Y\).
\end{enumerate}

This is the scheme-theoretic argument of Stacks Project, Tag 080A, rewritten
in the notation of the stated Lean theorem.

\end{document}
